\documentclass[11pt]{article}
\usepackage[margin=1in]{geometry}
\usepackage{amsmath,amssymb,amsthm}
\usepackage{booktabs}
\usepackage{array}
\usepackage[hidelinks]{hyperref}

\newtheorem{theorem}{Theorem}
\newtheorem{proposition}{Proposition}
\newtheorem{remark}{Remark}
\newcommand{\EqI}{\mathrel{\sim_I}}

\title{P3-T06 Draft/Control Audit and Stress:\\
The Unchanged Invariant-Functor EqSrc Candidate}
\author{Smuggling Auditor task overlay, RT-20260720-027}
\date{21 July 2026}

\begin{document}
\maketitle

\begin{center}
\fbox{\parbox{0.92\linewidth}{
\textbf{Authority and status.}
This is a \textbf{draft/control}, \textbf{proposal-only} audit containing
\textbf{source-extension data}.  The P3-T05 candidate is held byte-for-byte
fixed.  This packet neither repairs nor adopts nor rejects a source law.  It
does not establish physical observational equivalence, physical gauge,
general \(\mathrm{EqSrc}\), any downstream GR burden, proof authority,
publication authority, or a global no-go theorem.  The scoped status remains
\texttt{blocked\_adoption\_open\_continuation}.
}}
\end{center}

\section{Question, immutable input, and verdict}

The audited proposal defines
\[
 A\EqI B \quad\Longleftrightarrow\quad I(A)\cong I(B)
\]
for a prospectively fixed functor \(I:\mathcal C\to\mathcal D\).  Its
registered TeX SHA-256 is
\texttt{badc60c72ff16f84fe88568b825ff738a8f160a225f011715f027c1da2cfa1c3}.
The machine specification and finite witnesses are likewise held fixed at
\texttt{9c40dd2a83c5d9ee25a62c977b5b47eb5e2770890c8a7aa90dbe342cc7936ad1}
and
\texttt{6f101a2b3e7bc808e165349ab9cab33f44fcb671d2eb3f294ac8bf6071d38ecd}.

The decisive scoped verdict is
\begin{quote}
\texttt{conditional\_local\_source\_purity\_pass\_with\_arbitrary\_partition\_and\_interface\_completeness\_obstruction}.
\end{quote}
The conditional theorems are valid as written and no direct target or process
import was found.  The schema is nevertheless underdetermined as a source
relation: prospectivity alone permits every finite partition, and the fixed
homology witness forgets interface facts needed by a simple declared
operation.  The exact construction is therefore locally frozen as a generic
kernel schema pending independent source derivation and operation-completeness
evidence.  This is not a global rejection.

\section{Direct and transitive provenance audit}

The direct mathematical definition uses objects and morphisms of
\(\mathcal C\), a functor \(I\), and isomorphisms internal to
\(\mathcal D\).  The conditional operation results additionally use the
declared data \(\otimes,\boxtimes,\mu,W,L,Q,\bar I,V,\alpha\).  Desired
EqSrc labels, a target atlas, a target metric, benchmark outcomes, registry
rows, role outputs, validator results, file order, and process state are not
arguments to the relation or its proofs.  In this exact textual sense the
candidate is source-pure as written.

That pass is conditional and local.  Current canonical ontology derives none
of \(\mathcal C,\mathcal D,I\), the declared operation classes, or their
comparison maps.  The graph-homology instantiation is standard mathematics
applied to a declared source-extension model; it is not an ontology-derived
substrate invariant.  Registries establish identity and provenance, roles
bound permissions, and validators check syntax and finite computations.  None
of those operational surfaces supplies scientific truth or a missing source
law.  Absence of a direct target import therefore does not imply a source
derivation.

\section{Conditional theorem audit}

\begin{proposition}[Fixed theorem package]
\texttt{IFQ1}--\texttt{IFQ6} are correct under their stated hypotheses.
The marked-interface countermodel below does not refute any of them.
\end{proposition}
\begin{proof}
\texttt{IFQ1} is the pullback of isomorphism of objects along \(I\), so
identity, inverse, and composition of isomorphisms prove reflexivity,
symmetry, and transitivity.  In \texttt{IFQ2}, bifunctoriality and the
comparison isomorphisms \(\mu_{A,B}\) conjugate image isomorphisms into an
isomorphism after composition.  \texttt{IFQ3} and \texttt{IFQ4} correctly
separate preservation from reflection.  \texttt{IFQ5} is conjugation by the
components of an invertible natural transformation.  \texttt{IFQ6} correctly
notes that these categorical premises contain no physical observation,
matter-response, dynamical, or gauge semantics.
\end{proof}

This result certifies theorem validity, not candidate adequacy.  In
particular, \texttt{IFQ2} is a conditional congruence theorem; it is not a
claim that every intended source operation admits the required comparison.

\section{Finite arbitrary-partition theorem}

\begin{theorem}[Finite discrete partition realization]
Let \(X\) be a finite set and let \(\Pi\) be any partition of \(X\).  There
are finite discrete categories \(\mathcal C_X\) and \(\mathcal D_\Pi\) and a
functor \(I_\Pi:\mathcal C_X\to\mathcal D_\Pi\), fixed before any object
pair is evaluated, such that
\[
 x\sim_{I_\Pi}y \quad\Longleftrightarrow\quad
 x\text{ and }y\text{ lie in the same block of }\Pi.
\]
\end{theorem}
\begin{proof}
Take the objects of \(\mathcal C_X\) to be the elements of \(X\), with only
identity morphisms.  Take one object of the discrete category
\(\mathcal D_\Pi\) for each block of \(\Pi\).  Map each \(x\) to its block
and its identity to that block's identity.  This is a functor because there
are no nonidentity source arrows.  Two objects of a discrete target are
isomorphic exactly when they are equal, which holds exactly for elements of
the same block.
\end{proof}

For \(X=\{a,b,c\}\), the companion JSON enumerates all five partitions:
the discrete partition, the three two-block partitions, and the one-block
partition.  The first realizes the identity relation and the last realizes
the universal relation by a constant functor.  Hence prospective fixation is necessary to make
provenance auditable, but it is not a mathematical selection principle.  Any
finite desired partition can be fixed ``in advance'' unless an independent
source derivation, admissibility rule, nontriviality condition, or minimality
criterion restricts \(I\).

This theorem does not show that the P3-T05 text actually used desired labels;
no such use was found.  It shows that the stated prospectivity guard cannot by
itself establish that the relation is genuine rather than arbitrary.

\section{Fixed graph-homology stress}

The P3-T05 witness uses
\[
 H(G)=(\beta_0(G),\beta_1(G)),\qquad
 \beta_1=|E|-|V|+\beta_0.
\]
The following controls separate useful invariance from completeness:

\begin{center}
\begin{tabular}{>{\raggedright\arraybackslash}p{0.24\linewidth}
                >{\raggedright\arraybackslash}p{0.27\linewidth}
                >{\raggedright\arraybackslash}p{0.39\linewidth}}
\toprule
Stress & Exact result & Consequence \\
\midrule
Triangle vs. square & both \((1,1)\) & Relabel-independent homology is not
complete for graph isomorphism, vertex count, or edge count. \\
Edge vs. path vs. point & all \((1,0)\) & The quotient collapses connected
trees with different size and interface behavior. \\
Triangle relabeling & \((1,1)\) preserved & A pure source relabeling is
correctly ignored. \\
One edge vs. two edges & \((1,0)\ne(2,0)\) & Component count remains visible;
the invariant is not universal on the bounded domain. \\
\bottomrule
\end{tabular}
\end{center}

These are not defects in integral homology.  They delimit which source facts
the chosen object invariant can support.  A later source law may legitimately
ignore graph size, but that judgment requires source or physical semantics not
contained in the invariant itself.

\section{Marked-interface operation countermodel}

Consider finite simple graphs equipped with an ordered pair of distinct marked
vertices.  Let \(E\) add the simple edge between the marked endpoints when it
is absent and act as the identity when the edge is already present.  Let
\(X\) be \(K_2\) with its two vertices marked and let \(Y\) be the
three-vertex path \(P_3\) with its two degree-one endpoints marked.  Both are
connected trees, so
\[
 H(X)=H(Y)=(1,0), \qquad X\sim_HY.
\]
The marked vertices of \(X\) are already adjacent, so \(E(X)=X\) and
\(H(E(X))=(1,0)\).  The marked vertices of \(Y\) are not adjacent; adding
their edge creates a triangle, so \(H(E(Y))=(1,1)\).  Therefore
\[
 E(X)\not\sim_H E(Y).
\]

The H-kernel quotient has forgotten endpoint adjacency, an interface fact on
which \(E\) depends.  Thus this operation does not descend to that object
quotient.  This is a scoped countermodel to an arbitrary
composition-congruence, interface-completeness, or later-dynamics overread.
It is \emph{not} a refutation of \texttt{IFQ2}: no target-side operation and
no natural comparison isomorphism \(H(EG)\cong\bar E(HG)\) were supplied.
The countermodel proves failure of the missing premise in this domain.

\section{Layered disposition}

\begin{enumerate}
  \item \textbf{Theorem validity:} the conditional category-theoretic package
  survives unchanged.
  \item \textbf{Candidate adequacy:} the kernel schema is valid, but
  prospectivity alone permits arbitrary finite partitions, and the homology
  instantiation is not complete for marked-interface operations.
  \item \textbf{Source derivation:} current ontology derives none of the
  category, functor, or comparison package.  Adoption remains blocked while
  continuation remains open.
  \item \textbf{Physical meaning:} mathematical image isomorphism is not
  observational equivalence, dynamical interchangeability, or physical gauge.
\end{enumerate}

The exact P3-T05 construction is locally frozen as a generic kernel schema.
It may reopen only if a source-side derivation fixes \(\mathcal C,\mathcal
D,I\) without forbidden inputs, a nontriviality or minimality condition
excludes arbitrary partitions and identity/constant extremes, explicit
comparison data or a complete interface makes every intended operation
descend, and a separate physical bridge is supplied before physical
interpretation.  This scoped freeze preserves the candidate and its negative
evidence; it is neither an ontology rejection nor a theorem that future
source extensions are impossible.

The active milestone remains \texttt{source\_equivalence\_eqsrc}.  The burden
remains to determine whether a quotient candidate is a genuine source relation
rather than a disguised target or arbitrary partition.  No Distance-to-GR or
metric-use ledger row changes.  The dependency-ready next work item is
P3-T07, which may freeze superseded selector families, update the EqSrc theorem
inventory, and select one clean frontier.  P3-T07 is selected but not executed
here.

\end{document}
